\documentclass[11pt,a4paper]{article}
\usepackage[utf8]{inputenc}
\usepackage[T1]{fontenc}
\usepackage{lmodern}
\usepackage[portuguese]{babel}
\usepackage[a4paper,margin=2.5cm]{geometry}
\usepackage{xcolor}
\definecolor{TextEspresso}{HTML}{4A4238}
\definecolor{NudeTaupe}{HTML}{BFAFA6}
\definecolor{SoftBlush}{HTML}{D9C5B2}
\usepackage{titlesec}
\usepackage{enumitem}
\usepackage{listings}
\usepackage{hyperref}
\hypersetup{colorlinks=true, linkcolor=TextEspresso, urlcolor=TextEspresso}

\titleformat{\section}{\color{TextEspresso}\Large\bfseries}{\thesection}{0.5em}{}[{\color{NudeTaupe}\titlerule[0.8pt]}]
\titleformat{\subsection}{\color{TextEspresso}\large\bfseries}{\thesubsection}{0.5em}{}

\lstset{
  basicstyle=\ttfamily\small,
  breaklines=true,
  frame=single,
  rulecolor=\color{NudeTaupe},
  columns=fullflexible,
  showstringspaces=false
}

\newcommand{\guiaotitle}[2]{%
  \begin{center}
  {\Large\bfseries\color{TextEspresso} #1}\\[0.3cm]
  {\normalsize\color{NudeTaupe} #2}\\[0.2cm]
  {\color{NudeTaupe}\rule{4cm}{0.5pt}}
  \end{center}
  \vspace{0.5cm}
}

\begin{document}
\guiaotitle{Guião 05 -- Cluster Kubernetes em VMs Proxmox}{Infraestruturas de Computação na Nuvem, Aula 05}

\section{Objetivos}
\begin{itemize}[itemsep=0.2cm]
  \item Provisionar, a partir do Proxmox, as VMs que constituirão um cluster Kubernetes real com múltiplos nós.
  \item Inicializar o control plane com \texttt{kubeadm} e juntar um worker node ao cluster.
  \item Instalar um plugin de rede CNI e compreender o seu papel.
  \item Criar um Deployment e um Service a partir de manifestos YAML declarativos.
  \item Observar o agendamento de Pods entre nós e o mecanismo de self-healing.
\end{itemize}

\section{Pré-requisitos e Material}
\begin{itemize}[itemsep=0.2cm]
  \item Acesso ao Proxmox VE do servidor central e ao intervalo de VMIDs atribuído.
  \item Duas VMs Linux (Ubuntu Server LTS), criadas para este efeito: \texttt{k8s-cp} (control plane) e \texttt{k8s-w1} (worker). Recomenda-se 2 vCPU e 2~GB de RAM em cada, no mínimo.
  \item As duas VMs devem estar na mesma rede (mesma bridge do Proxmox) e alcançar-se mutuamente.
  \item Acesso SSH a ambas as VMs, com privilégios de \texttt{sudo}.
\end{itemize}

\textbf{Nota sobre recursos:} este guião cria duas VMs por aluno ou por grupo. Num servidor partilhado por uma turma inteira, isto pode esgotar a memória disponível. Confirme com a docente o modelo a adotar: por grupo (recomendado), ou individualmente com um número reduzido de nós. A alternativa mais leve consiste em cada grupo montar um cluster de dois nós, e partilhar observações entre grupos.

\section{Introdução}
Na Aula 04 corremos contentores num único host, com Docker. O passo seguinte -- e o tema desta aula -- é a orquestração: distribuir contentores por vários hosts, reagir automaticamente a falhas e expor serviços de forma estável.

Ao contrário de uma instalação de aprendizagem num único nó, aqui vamos construir um cluster \textbf{real, com dois nós}, sobre VMs provisionadas no Proxmox. Isto permite observar coisas que um cluster de nó único não mostra: o \emph{scheduler} a decidir em que nó colocar cada Pod, a rede entre nós assegurada por um plugin CNI, e a diferença entre control plane e worker node.

A ferramenta usada é o \texttt{kubeadm}, o instalador oficial do projeto Kubernetes para criar clusters conformes.

\section{Parte 1 -- Provisionar as VMs no Proxmox}
\begin{enumerate}[label=\arabic*.]

  \item \textbf{Criar as duas VMs.} A partir da interface web do Proxmox (ou por CLI, como no Guião 03), crie duas VMs Ubuntu Server. Se tiver criado um template no Guião 03, clone-o -- é bastante mais rápido:
\begin{lstlisting}
$ qm clone <id-template> <vmid-cp> --name k8s-cp --full
$ qm clone <id-template> <vmid-w1> --name k8s-w1 --full
$ qm set <vmid-cp> --memory 2048 --cores 2
$ qm set <vmid-w1> --memory 2048 --cores 2
$ qm start <vmid-cp> && qm start <vmid-w1>
\end{lstlisting}

  \item \textbf{Definir hostnames distintos.} Em cada VM, via SSH:
\begin{lstlisting}
# Na VM do control plane:
$ sudo hostnamectl set-hostname k8s-cp

# Na VM do worker:
$ sudo hostnamectl set-hostname k8s-w1
\end{lstlisting}
  Se as VMs foram clonadas do mesmo template, confirme também que têm endereços MAC e \texttt{machine-id} distintos (o Proxmox gera novos MACs ao clonar; para o \texttt{machine-id}, se necessário: \texttt{sudo rm -f /etc/machine-id \&\& sudo systemd-machine-id-setup}).

  \item \textbf{Registar os endereços IP} de ambas as VMs -- serão necessários adiante:
\begin{lstlisting}
$ ip addr show
\end{lstlisting}

\end{enumerate}

\section{Parte 2 -- Preparar ambos os nós}
Os passos seguintes executam-se \textbf{em ambas as VMs}.

\begin{enumerate}[label=\arabic*.,resume]

  \item \textbf{Desativar a swap.} O kubelet exige que a swap esteja desativada:
\begin{lstlisting}
$ sudo swapoff -a
$ sudo sed -i '/ swap / s/^/#/' /etc/fstab
\end{lstlisting}

  \item \textbf{Carregar os módulos de kernel necessários:}
\begin{lstlisting}
$ cat <<EOF | sudo tee /etc/modules-load.d/k8s.conf
overlay
br_netfilter
EOF
$ sudo modprobe overlay
$ sudo modprobe br_netfilter
\end{lstlisting}

  \item \textbf{Configurar os parâmetros de rede} para que o tráfego em bridge seja visto pelo iptables:
\begin{lstlisting}
$ cat <<EOF | sudo tee /etc/sysctl.d/k8s.conf
net.bridge.bridge-nf-call-iptables  = 1
net.bridge.bridge-nf-call-ip6tables = 1
net.ipv4.ip_forward                 = 1
EOF
$ sudo sysctl --system
\end{lstlisting}

  \item \textbf{Instalar o runtime de contentores (containerd).} Note que é o mesmo containerd referido na arquitetura Docker da Aula 04:
\begin{lstlisting}
$ sudo apt update
$ sudo apt install -y containerd
$ sudo mkdir -p /etc/containerd
$ containerd config default | sudo tee /etc/containerd/config.toml > /dev/null
\end{lstlisting}
  Ative o driver de cgroups \texttt{systemd}, requisito para o kubelet funcionar corretamente:
\begin{lstlisting}
$ sudo sed -i 's/SystemdCgroup = false/SystemdCgroup = true/' /etc/containerd/config.toml
$ sudo systemctl restart containerd
$ sudo systemctl enable containerd
\end{lstlisting}

  \item \textbf{Instalar \texttt{kubeadm}, \texttt{kubelet} e \texttt{kubectl}} a partir do repositório oficial. Confirme com a docente a versão a usar; o exemplo seguinte assume a série v1.30 (ajuste se necessário):
\begin{lstlisting}
$ sudo apt install -y apt-transport-https ca-certificates curl gpg
$ curl -fsSL https://pkgs.k8s.io/core:/stable:/v1.30/deb/Release.key \
    | sudo gpg --dearmor -o /etc/apt/keyrings/kubernetes-apt-keyring.gpg
$ echo 'deb [signed-by=/etc/apt/keyrings/kubernetes-apt-keyring.gpg] https://pkgs.k8s.io/core:/stable:/v1.30/deb/ /' \
    | sudo tee /etc/apt/sources.list.d/kubernetes.list
$ sudo apt update
$ sudo apt install -y kubelet kubeadm kubectl
$ sudo apt-mark hold kubelet kubeadm kubectl
\end{lstlisting}
  O \texttt{apt-mark hold} impede atualizações automáticas que poderiam quebrar o cluster.

\end{enumerate}

\section{Parte 3 -- Inicializar o cluster}
\begin{enumerate}[label=\arabic*.,resume]

  \item \textbf{Inicializar o control plane.} \emph{Apenas na VM \texttt{k8s-cp}}:
\begin{lstlisting}
$ sudo kubeadm init --pod-network-cidr=10.244.0.0/16
\end{lstlisting}
  O valor de \texttt{--pod-network-cidr} corresponde ao intervalo esperado pelo Flannel, o plugin CNI que instalaremos a seguir. \textbf{Guarde o comando \texttt{kubeadm join} apresentado no final} -- será necessário no passo seguinte.

  \item \textbf{Configurar o \texttt{kubectl}} para o utilizador normal, ainda em \texttt{k8s-cp}:
\begin{lstlisting}
$ mkdir -p $HOME/.kube
$ sudo cp -i /etc/kubernetes/admin.conf $HOME/.kube/config
$ sudo chown $(id -u):$(id -g) $HOME/.kube/config
$ kubectl get nodes
\end{lstlisting}
  O nó aparecerá em estado \texttt{NotReady} -- falta a rede.

  \item \textbf{Observar porque o nó não está pronto.} Confirme que os Pods de sistema relacionados com DNS estão pendentes:
\begin{lstlisting}
$ kubectl get pods -n kube-system
\end{lstlisting}
  Os Pods do CoreDNS ficam em \texttt{Pending} porque ainda não existe rede de Pods. Isto demonstra na prática o papel do CNI discutido na teoria.

  \item \textbf{Instalar o plugin CNI (Flannel):}
\begin{lstlisting}
$ kubectl apply -f https://github.com/flannel-io/flannel/releases/latest/download/kube-flannel.yml
\end{lstlisting}
  Aguarde alguns instantes e volte a verificar:
\begin{lstlisting}
$ kubectl get nodes
$ kubectl get pods -n kube-system
\end{lstlisting}
  O nó deve agora passar a \texttt{Ready} e o CoreDNS a \texttt{Running}.

  \item \textbf{Juntar o worker ao cluster.} \emph{Na VM \texttt{k8s-w1}}, execute o comando \texttt{kubeadm join} guardado anteriormente:
\begin{lstlisting}
$ sudo kubeadm join <IP-do-cp>:6443 --token <token> \
    --discovery-token-ca-cert-hash sha256:<hash>
\end{lstlisting}
  Se tiver perdido o comando, pode gerá-lo de novo a partir do control plane:
\begin{lstlisting}
$ kubeadm token create --print-join-command
\end{lstlisting}

  \item \textbf{Confirmar o cluster completo,} a partir de \texttt{k8s-cp}:
\begin{lstlisting}
$ kubectl get nodes -o wide
\end{lstlisting}
  Devem aparecer dois nós em estado \texttt{Ready}, com papéis (\texttt{ROLES}) distintos.

  \item \textbf{Inspecionar os componentes do control plane.} Identifique os componentes estudados na teoria a correr como Pods:
\begin{lstlisting}
$ kubectl get pods -n kube-system -o wide
\end{lstlisting}
  Localize o \texttt{kube-apiserver}, o \texttt{etcd}, o \texttt{kube-scheduler} e o \texttt{kube-controller-manager}, e repare em que nó cada um corre.

\end{enumerate}

\section{Parte 4 -- Implantar uma aplicação}
\begin{enumerate}[label=\arabic*.,resume]

  \item \textbf{Criar o manifesto do Deployment}, \texttt{deployment.yaml}, com pedidos e limites de recursos:
\begin{lstlisting}
apiVersion: apps/v1
kind: Deployment
metadata:
  name: web-app
  labels:
    app: web-app
spec:
  replicas: 4
  selector:
    matchLabels:
      app: web-app
  template:
    metadata:
      labels:
        app: web-app
    spec:
      containers:
        - name: web-app
          image: nginx:1.25
          ports:
            - containerPort: 80
          resources:
            requests:
              memory: "64Mi"
              cpu: "100m"
            limits:
              memory: "128Mi"
              cpu: "250m"
          readinessProbe:
            httpGet:
              path: /
              port: 80
            initialDelaySeconds: 3
            periodSeconds: 5
\end{lstlisting}

  \item \textbf{Aplicar o manifesto e observar a distribuição pelos nós:}
\begin{lstlisting}
$ kubectl apply -f deployment.yaml
$ kubectl get deployments
$ kubectl get replicasets
$ kubectl get pods -o wide
\end{lstlisting}
  A coluna \texttt{NODE} mostra em que nó o scheduler colocou cada Pod. Registe a distribuição obtida.

  \item \textbf{Criar o Service}, \texttt{service.yaml}:
\begin{lstlisting}
apiVersion: v1
kind: Service
metadata:
  name: web-app-service
spec:
  selector:
    app: web-app
  ports:
    - port: 80
      targetPort: 80
  type: NodePort
\end{lstlisting}

  \item \textbf{Aplicar o Service e aceder à aplicação:}
\begin{lstlisting}
$ kubectl apply -f service.yaml
$ kubectl get services
\end{lstlisting}
  Identifique a porta atribuída (no intervalo 30000--32767) e aceda através de \emph{qualquer} um dos nós:
\begin{lstlisting}
$ curl http://<IP-do-k8s-cp>:<nodeport>
$ curl http://<IP-do-k8s-w1>:<nodeport>
\end{lstlisting}
  Confirme que ambos respondem -- é o \texttt{kube-proxy} a encaminhar o tráfego, independentemente do nó onde os Pods realmente estão.

  \item \textbf{Testar a resolução de nomes interna (CoreDNS).} A partir de um Pod temporário dentro do cluster:
\begin{lstlisting}
$ kubectl run teste --rm -it --image=busybox:1.36 --restart=Never -- sh
/ # wget -O- http://web-app-service
/ # exit
\end{lstlisting}

\end{enumerate}

\section{Parte 5 -- Escalar e observar o self-healing}
\begin{enumerate}[label=\arabic*.,resume]

  \item \textbf{Escalar o número de réplicas} e observar onde são colocadas as novas:
\begin{lstlisting}
$ kubectl scale deployment web-app --replicas=6
$ kubectl get pods -o wide
\end{lstlisting}

  \item \textbf{Apagar um Pod manualmente e observar a reconciliação:}
\begin{lstlisting}
$ kubectl get pods -l app=web-app
$ kubectl delete pod <nome-de-um-pod>
$ kubectl get pods -l app=web-app --watch
\end{lstlisting}
  Observe que, poucos segundos depois, o ReplicaSet cria um novo Pod. Interrompa com \texttt{Ctrl+C} quando confirmar.

  \item \textbf{Simular a falha de um nó inteiro.} Este é o teste que um cluster de nó único não permite fazer. A partir do Proxmox, pare abruptamente a VM do worker:
\begin{lstlisting}
$ qm stop <vmid-w1>
\end{lstlisting}
  A partir de \texttt{k8s-cp}, observe a reação do cluster ao longo de alguns minutos:
\begin{lstlisting}
$ kubectl get nodes --watch
$ kubectl get pods -o wide
\end{lstlisting}
  O nó passa a \texttt{NotReady} e, decorrido o tempo de tolerância configurado, os Pods que lá corriam são recriados no nó restante. Registe quanto tempo demorou até observar a recriação.

  \item \textbf{Recuperar o nó} e confirmar que volta a integrar o cluster:
\begin{lstlisting}
$ qm start <vmid-w1>
$ kubectl get nodes
\end{lstlisting}

  \item \textbf{Limpeza.} Remova os recursos aplicados. Mantenha o cluster montado se for reutilizado nas Aulas 08 e 09 (confirmar com a docente); caso contrário, pare também as VMs:
\begin{lstlisting}
$ kubectl delete -f service.yaml
$ kubectl delete -f deployment.yaml
\end{lstlisting}

\end{enumerate}

\section{Entregável / Questões de Verificação}
\begin{enumerate}[label=\alph*)]
  \setlength\itemsep{0.2cm}
  \item Submeta os ficheiros \texttt{deployment.yaml} e \texttt{service.yaml}, e a saída de \texttt{kubectl get nodes -o wide} com os dois nós em \texttt{Ready}.
  \item Antes de instalar o Flannel, o nó estava \texttt{NotReady} e o CoreDNS em \texttt{Pending}. Explique porquê, com base no papel do plugin CNI no modelo de rede do Kubernetes.
  \item Registe a distribuição de Pods pelos nós (\texttt{kubectl get pods -o wide}). Que componente do control plane tomou essa decisão e que informação usou?
  \item Acedeu à aplicação através do IP de ambos os nós, mesmo não havendo Pods em todos eles. Que componente torna isto possível?
  \item Quanto tempo decorreu entre parar a VM do worker e a recriação dos seus Pods no outro nó? Relacione o comportamento observado com o ciclo de reconciliação estudado na teoria.
  \item Qual a diferença de responsabilidades entre o Deployment e o ReplicaSet que ele cria?
  \item Definiu \texttt{requests} e \texttt{limits} no manifesto. Qual a diferença entre os dois, e através de que mecanismo do kernel (visto na Aula 04) são os limites efetivamente aplicados?
  \item Este cluster corre sobre VMs que, por sua vez, correm sobre o Proxmox. Que vantagem tem esta arquitetura em relação a instalar o Kubernetes diretamente no servidor físico?
\end{enumerate}

\end{document}
